\section{I had too much to Dream last Night}

I had a dream last night that has left me bemused. Without going into what happened before, most of which I have immediately forgotten, the dream ended with a brief conversation with a friend at a café. I was explaining to her that she had two ways to look at the world.

\begin{enumerate}
\item Everything that happens depends on the past. The effects of the mass of past events weighs heavy on our current options, limiting our field of action, so it is difficult to extricate ourselves from what has gone on before. As we grow older, we sense this more and more. As the human race grows older, as a race we have fewer options.

\item Each moment is brand new, independent of what may have happened before. Hence, at any instant a new world is possible.
\end{enumerate}


I've heard stories of people receiving answers in their dreams. I just get a dilemma.

\url{https://www.youtube.com/watch?v=m7KxShixE-A}

\flrightit{Posted on 2011-11-24 by Cologero}

\begin{center}* * *\end{center}

\begin{footnotesize}
\begin{sffamily}
\texttt{Jupiter on 2011-11-24 at 14:50 said:}

What dilemma? Neither of those possibilities can be true without the simultaneous actuality of the other possibility.

The duality is solely in your mind.

\hfill

\texttt{Logres on 2011-11-24 at 19:25 said:}

But doesn't it have to be ``in'' your mind first, before you can transcend it? Unless someone is passing it by touch, or you've ``done it before somewhere else'', but then, perhaps the same thing? And could someone transcend it by having a dream, and then distancing one's self from the dilemma that way?

\hfill

\texttt{charlotte on 2011-11-25 at 12:38 said:}

it's said that mundane man is controlled by his environment, while superior man governs his environment. There is a divine plan to all creation and existence, but the free will of man is a constantly limiting factor. Regarding `all potential' being inherent in each moment, Timothy Leary opined that the universe is moving towards a state of ever-incresing complexity with the goal of a state of complete novelty. By achieving novelties within the infinite potential we gradually attain increased connectivity with the rest of the universe -- by joining up the dots so to speak. A moment will then come when connectivity is finally reached, at which point there will be harmony because of natural understanding between beings. It's a matter of time, basically... .

\hfill

\texttt{charlotte on 2011-11-25 at 12:40 said:}

i think you can transcend a former state without having prior consciousness of it

\hfill

\texttt{logres on 2011-11-25 at 21:12 said:}

Your comments, long as they often are, always have a profitable point. You've helped connect the dots.

\hfill

\texttt{Perennial on 2011-12-02 at 02:41 said:}

I have never fully grasped the traditional view of reincarnation, but as I understand it, if we are understanding the world through the eyes of a spirit which has existed before, could we not, by intuition, know of matters without ``knowing'' them, but harvest from the knowledge of our soul? Would this not help our soul be more advanced when we entered the next realm? Surely understanding the world through the spirit is easier then always trying to grasp it mentally in the physical plane only.

\hfill

\texttt{Perennial on 2011-12-02 at 02:41 said:}

I doubt anything I just said made a bit of sense.

\hfill

\texttt{Charlotte Cowell on 2011-12-02 at 14:40 said:}

yes it did make sense. In answer to your questions, reincarnation is something which (by general though not unanimous consensus) either happens or has happened. Whether it is (now) necessary is more ponderable. Whether conscious recollection of past lives is either necessary and/or beneficial for spiritual growth is highly debatable. I would say not in the vast majority of cases, although there may well be unconscious and subconscious benefits. When it comes to remembrance of (possible) former lives, I'd say it can be hazardous, depending on how the knowledged is attained, and may only happen in either emergency situations or cases where an individual has such a high degree of spiritual attainment that they have the potential to achieve a Boddhisattava type state and work great healing in the world. (It stands to reason that the enemy might also make use of such insighsts). Generally, though, I think it is better not to see that much and ultimately reincarnation is an unnecessary doctrine because time and space have no meaning in the spiritual dimension. We are all Adam or Eve (both, even, if the universal androgyne is considered!), what more do we need to know?

In Letter V of meditations on the Tarot (The Pope), the author speaks of reincarnation:

The Buddha recognised and at the same time denied the fact of reincarnation. He recognised it as fact and he denied it as ideal. Because facts are transitiory; they come and go.

There was a time when there was no reincarnation; there will be a time when it will no longer be. Reincarnation commenced only after the Fall and it will cease with Reintegration. It is therefore not eternal, and therefore it is not an ideal.

A longer extract is here:

\url{http://alchemical-weddings.com/alchemical-weddings/two-truths}


On the other hand pertaining more to your comment, certain polarised archetypes -- of the kind that appear most often in reincarnation memories (or `memories') -- can certainly be benefit to contemplate. Clearly if one sees one's self as a reincarnation of Apollo, Caesar or John the Baptist there is much to be learned from that particular blueprint of the soul! Similarly, understanding what it is to be miserly, spiteful and therefore lonely -- as if one once was in another, more ordinary life -- might well give one the desire to be more generous, cheerful and sociable. But in either extreme reincarnation is necessary neither in fact nor in theory. Empathy will suffice and compassion will bring the wisdom to use the knowledge in a beneficial way.

Empathy in itself helps to generate that compassion, and likely also a state of heightened perception, intuition, joy and other valuable spiritual fruits. It can also be very painful. As can a knowledge of reincarnation -- as in, knowing (or `knowing') of one's past lives. Some schools of thought woud say that in one's `last life' this type of memory occurs, although this might also be seen as a consequence of initiation, which permits conscious identification with and connectivity to the whole, with Divinity, unfettered by time or space. In either case recollection of one's past life dramas -- be they real or clever mental constructs -- may well be used to deal with the karma of the soul, but this is an exceedingly complex procedure involving participation from one's karmic relations and friends, ideally with some degree of consciousness also, though perhaps not of the identical pool of knowledge. Each to their own after all.

Either way, in one or many lives, one must pick up and put back together the pieces of the self in the process of returning to source. Connect the dots!

\hfill

\texttt{Charlotte Cowell on 2011-12-03 at 06:58 said:}

the other thing I was going to say about that but forgot initially, is to do with the lake of forgettfulness, as it strikes me that one may have to drink from this if one has previously been awakened to reincarnation memories, from which it may be almost impossible to move on. By which I mean, the pull of one or more former lives might be stronger than the desire to lead one's present life, thus hindering development. And even Buddha said that to forgive is good but to forget is better... . but in this case, one must then have a `backup plan' with respect to the divine/ blueprint -- perhaps this is one function of the hermit's `book'


\hfill

\end{sffamily}
\end{footnotesize}

